\documentclass[11pt]{article}
\usepackage[margin=1in]{geometry}
\usepackage{amsmath,amssymb,bm}
\usepackage{booktabs}
\usepackage{enumitem}
\usepackage[hidelinks]{hyperref}

\newcommand{\code}[1]{\texttt{#1}}
\newcommand{\Norm}{\mathcal{N}}
\newcommand{\pstay}{p_{\mathrm{stay}}}
\newcommand{\x}{\mathbf{x}}

\title{Bin-cell HMM for CHISEL:\\model, fitting and copy-number calling}
\author{Runpeng Luo}
\date{2026-11-12}

\begin{document}
\maketitle

\begin{abstract}
CHISEL~\cite{zaccaria2021} clusters genomic bins globally: each bin is one point of RDRs and BAFs across all cells,
so a cluster is a set of bins shared by every cell. The bin-cell HMM clusters the
(bin, cell) points instead. Every cell is a Markov chain along its bins over $K$ hidden states
shared by all cells, each state emitting (RDR, BAF) from a Gaussian. The emission means start
on the integer copy-number grid and are refined by Baum-Welch with fixed start and transition
probabilities. The states are then labeled with copy numbers, either by the nearest grid point
or, per cell, by CHISEL's own \code{Caller.calling}.
\end{abstract}

\section{Data}\label{sec:data}

Cells are $e = 1, \dots, N$, sorted by name, and bins $t = 1, \dots, T$, sorted by chromosome
and position. The bins of chromosome $c$ form the contiguous block
$\mathcal{T}_c = \{t_c, \dots, t'_c\}$.

For bin $t$ and cell $e$, CHISEL's \code{combo.tsv} holds the read count of the matched normal
$\nu_t$, the read count of the cell $n_{e,t}$, the RDR $r_{e,t}$ and the allele counts
$A_{e,t}$, $B_{e,t}$, with
\begin{align}
r_{e,t} &= \frac{n_{e,t}}{\nu_t}\cdot\frac{\sum_{t'}\nu_{t'}}{\sum_{t'} n_{e,t'}}
  \quad (0 \text{ if } \nu_t = 0), \label{eq:rdr}\\
\beta_{e,t} &= \frac{B_{e,t}}{A_{e,t}+B_{e,t}}
  \quad \bigl(\tfrac12 \text{ if } A_{e,t}+B_{e,t} = 0\bigr), \qquad
\tilde\beta_{e,t} = \min(\beta_{e,t},\, 1-\beta_{e,t}). \label{eq:baf}
\end{align}
The observation is $\x_{e,t} = (x_{e,t,1}, x_{e,t,2}) = (r_{e,t}, \beta_{e,t})$ with the phased
BAF on $[0, 1]$: B is taken as the same haplotype in every bin and cell, and $\beta_{e,t}$ is not
mirrored. Phase switches between bins are not modeled; a switch reads as a change of state. A value
that is not finite is missing: $m_{e,t,d} = 1$ if $x_{e,t,d}$ is observed and $0$ otherwise; an
entry without SNP-covering reads has no BAF.

CHISEL's \code{Caller.format\_combo} arranges the data as
$\mathbf{D} \in \mathbb{R}^{T \times 2N}$ with $D_{t,2e-1} = r_{e,t}$ and
$D_{t,2e} = \tilde\beta_{e,t}$, the mirrored BAF. \code{Clusterizer.kclustering} clusters the
$T$ rows of $\mathbf{D}$, points of $\mathbb{R}^{2N}$. \code{hmm\_caller.format\_phased} builds
the same layout with $D_{t,2e} = \beta_{e,t}$, and the bin-cell HMM models its $NT$ points
$\x_{e,t} \in \mathbb{R}^2$.

\section{Model}\label{sec:model}

The parameters are the number of states $K$, the fixed self-transition probability $\pstay$,
and the emission parameters $\theta = \{\bm\mu_k, \bm\sigma^2_k\}_{k=1}^{K}$ with
$\bm\mu_k, \bm\sigma^2_k \in \mathbb{R}^2$.

\subsection{Hidden states}\label{sec:states}

Each cell $e$ has a hidden state $z_{e,t} \in \{1, \dots, K\}$ at every bin. The $K$ states,
and so their emission parameters, are shared by all cells. Given $\theta$, the chains of
different cells, and of different chromosomes of one cell, are independent.

\subsection{Start and transition probabilities}\label{sec:transitions}

For every cell $e$, chromosome $c$ and bins $t-1, t \in \mathcal{T}_c$,
\begin{equation}
P(z_{e,t_c} = k) = \pi_k = \frac1K, \qquad
P(z_{e,t} = l \mid z_{e,t-1} = k) = A_{kl} =
\begin{cases}
\pstay & k = l,\\[2pt]
\dfrac{1-\pstay}{K-1} & k \ne l,
\end{cases}
\label{eq:transitions}
\end{equation}
and $\mathbf{A} = [1]$ when $K = 1$. $\bm\pi$ and $\mathbf{A}$ are fixed a priori, shared by
all cells, and not learned. $\mathbf{A}$ is doubly stochastic, and every entry is positive when
$0 < \pstay < 1$.

\subsection{Emission}\label{sec:emission}

\begin{equation}
p(\x_{e,t} \mid z_{e,t} = k) = b_k(\x_{e,t})
= \prod_{d=1}^{2} \Norm\bigl(x_{e,t,d};\, \mu_{k,d},\, \sigma^2_{k,d}\bigr)^{m_{e,t,d}} .
\label{eq:emission}
\end{equation}
A missing dimension contributes the factor 1, i.e., it is marginalized out.

\subsection{Likelihood}\label{sec:likelihood}

\begin{equation}
\mathcal{L}(\theta) = \prod_{e=1}^{N}\prod_{c}\ \sum_{z_{e,t_c},\dots,z_{e,t'_c}}
\pi_{z_{e,t_c}}\, b_{z_{e,t_c}}(\x_{e,t_c})
\prod_{t=t_c+1}^{t'_c} A_{z_{e,t-1} z_{e,t}}\, b_{z_{e,t}}(\x_{e,t}) .
\label{eq:likelihood}
\end{equation}

\section{Initialization on the copy-number grid}\label{sec:init}

\subsection{Grid}\label{sec:grid}

With maximum total copy number $C_{\max}$ and base $\rho$, the total copy number whose expected
RDR is 1,
\begin{equation}
\mathcal{G} = \bigl\{(a,b) \in \mathbb{N}_0^2 : a+b \le C_{\max}\bigr\},
\qquad
\mathbf{g}(a,b) = \Bigl(\frac{a+b}{\rho},\ \frac{b}{a+b}\Bigr), \quad
\mathbf{g}(0,0) = \bigl(0, \tfrac12\bigr).
\label{eq:grid}
\end{equation}
The grid holds both $(a, b)$ and $(b, a)$, at BAF $b/(a+b)$ and $a/(a+b)$: no emission mean is
constrained to BAF $\le \tfrac12$.

\subsection{Constrained k-means}\label{sec:kmeans}

Let $\mathcal{P}$ be the points $\x_{e,t}$ with both values observed, or a uniform subsample of
$M$ of them. Constrained k-means chooses $K$ distinct grid points
$\mathcal{S} = \{s_1, \dots, s_K\} \subset \mathcal{G}$ that minimize
\begin{equation}
J(\mathcal{S}) = \sum_{\x \in \mathcal{P}} \min_{k}\ \lVert \x - \mathbf{g}(s_k) \rVert^2 ,
\label{eq:kmeans}
\end{equation}
by Lloyd iterations~\cite{lloyd1982,isl2021,mackay2003} whose centers are restricted to the grid:
\begin{enumerate}[nosep]
\item assign each point to its nearest center,
  $\ell(\x) = \arg\min_k \lVert \x - \mathbf{g}(s_k) \rVert^2$;
\item in decreasing order of cluster size, move $s_k$ to the grid point nearest to the mean of
  its members that no larger cluster has taken (an empty cluster targets its current point);
\item stop when no center moves, or after \code{maxiter} iterations.
\end{enumerate}
Restart 0 starts from the $K$ grid points that are the nearest grid point of the most points.
Restart $r > 0$ draws $K$ grid points without replacement, each with probability proportional to
that count plus $10^{-9}$. The restart with the smallest $J$ is kept.

\subsection{Initial parameters}\label{sec:init-params}

\begin{equation}
\bm\mu_k = \mathbf{g}(s_k), \qquad
\sigma^2_{k,d} = \max\Bigl(\epsilon,\ \frac{1}{|\mathcal{P}|}\sum_{\x \in \mathcal{P}}
  \bigl(x_d - \mu_{\ell(\x),d}\bigr)^2\Bigr),
\label{eq:init}
\end{equation}
a variance tied across states, with floor $\epsilon$.

\section{Baum-Welch on the emission parameters}\label{sec:em}

EM maximizes $\log\mathcal{L}(\theta)$ over $\theta$ with $\bm\pi$ and $\mathbf{A}$ fixed.

\subsection{E-step: scaled forward-backward}\label{sec:estep}

For one cell and one chromosome block, index the bins $t = 1, \dots, \tau$ and drop $e$. Let
$h_t = \max_k \log b_k(\x_t)$ and $\tilde b_t(k) = \exp\bigl(\log b_k(\x_t) - h_t\bigr) \le 1$.
The scaled recursions~\cite{rabiner1989,mackay2003} are
\begin{align}
\hat\alpha_1(k) &= \frac{\pi_k\, \tilde b_1(k)}{c_1}, &
\hat\alpha_t(l) &= \frac{\tilde b_t(l)}{c_t} \sum_{k} \hat\alpha_{t-1}(k)\, A_{kl}, \label{eq:forward}\\
\hat\beta_\tau(k) &= 1, &
\hat\beta_t(k) &= \frac{1}{c_{t+1}} \sum_{l} A_{kl}\, \tilde b_{t+1}(l)\, \hat\beta_{t+1}(l), \label{eq:backward}
\end{align}
where $c_t$ normalizes $\hat\alpha_t$ to sum to 1. The posteriors and the log-likelihood are
\begin{equation}
\gamma_t(k) = P(z_t = k \mid \x_{1:\tau})
= \frac{\hat\alpha_t(k)\,\hat\beta_t(k)}{\sum_{k'} \hat\alpha_t(k')\,\hat\beta_t(k')},
\qquad
\log p(\x_{1:\tau}) = \sum_{t=1}^{\tau} \log c_t + \sum_{t=1}^{\tau} h_t .
\label{eq:posterior}
\end{equation}
Since $\tilde b_t(k) = 1$ for some $k$ and $\hat\alpha_{t-1}$ sums to 1,
$c_t \ge \min_{k,l} A_{kl} > 0$ and $c_1 \ge 1/K$, so the recursions do not underflow. All cells
of a chromosome block are processed together, vectorized over $e$.

\subsection{M-step}\label{sec:mstep}

Because $\bm\pi$ and $\mathbf{A}$ are fixed, the expected complete-data log-likelihood depends
on $\theta$ only through the emission terms, and the pairwise posteriors
$P(z_{t-1}, z_t \mid \x)$ are not needed:
\begin{equation}
Q(\theta) = \sum_{e,t,k} \gamma_{e,t}(k) \sum_{d} m_{e,t,d}\,
\log \Norm\bigl(x_{e,t,d};\, \mu_{k,d},\, \sigma^2_{k,d}\bigr) + \text{const}.
\end{equation}
With the sufficient statistics
\begin{equation}
S^{(j)}_{k,d} = \sum_{e,t} \gamma_{e,t}(k)\, m_{e,t,d}\, x_{e,t,d}^{\,j}, \qquad j = 0, 1, 2,
\end{equation}
the update is
\begin{equation}
\mu_{k,d} = \frac{S^{(1)}_{k,d}}{S^{(0)}_{k,d}}, \qquad
\sigma^2_{k,d} = \max\Bigl(\epsilon,\ \frac{S^{(2)}_{k,d}}{S^{(0)}_{k,d}} - \mu_{k,d}^2\Bigr),
\label{eq:mstep}
\end{equation}
and $\mu_{k,d}$, $\sigma^2_{k,d}$ keep their values where $S^{(0)}_{k,d} \le 10^{-8}$. For fixed
$\mu_{k,d}$, $Q$ is unimodal in $\sigma^2_{k,d}$, so \eqref{eq:mstep} maximizes $Q$ over
$\sigma^2 \ge \epsilon$. Each iteration is thus an EM step on the constrained parameter space,
and $\mathcal{L}$ does not decrease~\cite{dempster1977}.

\subsection{Convergence}\label{sec:convergence}

The E-step of iteration $i$ yields $\ell^{(i)} = \log\mathcal{L}(\theta^{(i)})$. Fitting stops at
the first $i \ge 1$ with $\ell^{(i)} - \ell^{(i-1)} < \text{tol}$, keeping $\theta^{(i)}$, or
after \code{maxiter} E-steps.

\section{Decoding}\label{sec:decoding}

For each cell and chromosome block, the Viterbi path~\cite{viterbi1967} is
\begin{align}
\delta_1(k) &= \log\pi_k + \log b_k(\x_1), &
\delta_t(l) &= \max_k \bigl[\delta_{t-1}(k) + \log A_{kl}\bigr] + \log b_l(\x_t),\\
\psi_t(l) &= \arg\max_k \bigl[\delta_{t-1}(k) + \log A_{kl}\bigr], &
\hat z_\tau &= \arg\max_k \delta_\tau(k), \quad \hat z_{t-1} = \psi_t(\hat z_t).
\end{align}
\code{hmmclustering} returns \code{members[t, e]} $= \hat z_{e,t}$. The posteriors
$\gamma_{e,t}(k)$ of \eqref{eq:posterior} are available from \code{predict\_proba}.

\section{Copy-number labels}\label{sec:labels}

\subsection{Grid labels}\label{sec:grid-labels}

State $k$ is labeled with the grid point nearest to its fitted mean,
\begin{equation}
(a_k, b_k) = \arg\min_{(a,b) \in \mathcal{G}} \lVert \bm\mu_k - \mathbf{g}(a,b) \rVert^2 ,
\end{equation}
the same label in every cell. This assumes that every cell has the RDR scale of
\eqref{eq:grid}, i.e., mean total copy number $\rho$.

\subsection{Per-cell labels with CHISEL's calling}\label{sec:calling}

\code{hmm\_caller.py} labels the states per cell with CHISEL's \code{Caller.calling}, unchanged.
For cell $e$, the cluster of state $k$ is $C_{k,e} = \{t : \hat z_{e,t} = k\}$, summarized as
\code{Caller.est\_clusters} summarizes a CHISEL cluster:
\begin{equation}
R_{k,e} = \frac{\sum_{t \in C_{k,e}} n_{e,t}\, r_{e,t}}{\sum_{t \in C_{k,e}} n_{e,t}}, \qquad
\tilde B_{k,e} = \frac{\sum_{t \in C_{k,e}} (A_{e,t}+B_{e,t})\, \tilde\beta_{e,t}}
                      {\sum_{t \in C_{k,e}} (A_{e,t}+B_{e,t})}, \qquad
\mathrm{size}_{k,e} = \sum_{t \in C_{k,e}} \mathrm{len}(t),
\end{equation}
with $R_{k,e} = 0$ and $\tilde B_{k,e} = \tfrac12$ where the denominators are 0. The steps
below drop $e$. Let $f(A,B) = \min(A,B)/(A+B)$, the mirrored BAF of a state.

\paragraph{Base cluster (\code{identify\_base}).}
Starting from $s = \delta$ (shift) and raising $s$ by $\delta/2$ up to $\tfrac12$, the balanced
clusters are
\begin{equation}
\mathcal{N}_s = \bigl\{k : |\tfrac12 - \tilde B_k| \le 0.2 + s,\ R_k > 0,\ |\tfrac12 - \hat\beta_k| \le s\bigr\},
\end{equation}
where $\hat\beta_k$ is the best of 50 EM runs fitting the mixture
$x_t \sim \tfrac12\,\mathrm{Bin}(A_t+B_t, \beta) + \tfrac12\,\mathrm{Bin}(A_t+B_t, 1-\beta)$
to the bins of $C_k$, $x_t$ being $A_t$ or $B_t$ at random (\code{Combiner.EM}). With
\code{--fastscaling} the condition is $|\tfrac12 - \tilde B_k| \le s$ and $R_k > 0$. At the
first $s$ with $\mathcal{N}_s \ne \emptyset$, every $k \in \mathcal{N}_s$ defines the bubble
$U_k = \{k' \in \mathcal{N}_s : |R_{k'} - R_k| \le 0.05\}$, and the base RDR $\bar R$ is the
size-weighted mean RDR of the bubble with the largest total size.

\paragraph{Scale.}
For every even base ploidy $p \in \{2, 4, \dots, P\}$, $\lambda_p = p / \bar R$: a state $(A,B)$
predicts RDR $(A+B)/\lambda_p$ and mirrored BAF $f(A,B)$.

\paragraph{Variance floors.}
For each $p$, cluster $k$ gets the provisional state $\tilde s_k$ with total
$\max(1, \mathrm{round}(R_k \lambda_p))$ whose $f$ is nearest to $\tilde B_k$. For each state
$s$, $V^B_p(s)$ and $V^R_p(s)$ are the sample variances of the bins' $\tilde\beta_t$ around
$f(s)$ and of their $r_t$ around $(A_s+B_s)/\lambda_p$, over the bins of the clusters with
$\tilde s_k = s$ (denominator $n-1$, or 1 if $n \le 1$; $0.001$ if $n = 0$). The floors are
$\bar v^B = \max_p \operatorname{mean}_s V^B_p(s)$ and
$\bar v^R = \max_p \operatorname{mean}_s V^R_p(s)$, the mean taken over every state with total
$1, \dots, \max_k \mathrm{round}(R_k\lambda_p)$.

\paragraph{State of each cluster.}
With $\hat\tau_k = \mathrm{round}(R_k \lambda_p)$, the candidates are the $(A,B)$ with $A \ge B$
and $A + B \in \{\hat\tau_k - 1, \hat\tau_k, \hat\tau_k + 1\} \cap [0, \max_{k'}\hat\tau_{k'}]$,
total 0 meaning $(0,0)$. The cluster takes the candidate maximizing
\begin{equation}
\ell_k(A,B) = \sum_{t \in C_k} \Bigl[\log\Norm\bigl(\tilde\beta_t;\, f(A,B),\, v^B_k\bigr)
  + \log\Norm\bigl(r_t;\, (A+B)/\lambda_p,\, v^R_k\bigr)\Bigr],
\end{equation}
with $v^B_k = \max\bigl(\bar v^B, \widehat{\mathrm{Var}}_{t \in C_k}(\tilde\beta_t \mid f(A,B))\bigr)$
and $v^R_k = \max\bigl(\bar v^R, \widehat{\mathrm{Var}}_{t \in C_k}(r_t \mid (A+B)/\lambda_p)\bigr)$,
the sample variances around the predicted values. A cluster with $R_k = 0$ takes $(0,0)$ with
score 1.

\paragraph{Base ploidy (BIC).}
Let $\tau^*_p$ be the largest total copy number whose clusters cover at least $2\%$ of the size
of the clusters with $R_k > 0$ (if none, $\max_k \hat\tau_k$), and
\begin{equation}
|\mathcal{S}_p| = 1 + \sum_{\tau=0}^{\tau^*_p} N(\tau), \qquad N(0) = 1,\ N(\tau) = \lfloor \tau/2 \rfloor + 1,
\end{equation}
the number of allele-specific states with total at most $\tau^*_p$, plus one. With $n$ the number
of bins of the cell and sensitivity $\kappa$ (\code{-A}),
\begin{equation}
\mathrm{BIC}(p) = 2\kappa\, |\mathcal{S}_p| \log n - 2 \sum_k \ell_k(A_k, B_k),
\label{eq:bic}
\end{equation}
and the cell takes the $p$ of smallest BIC among those whose clusters assigned $(0,0)$ cover less
than the fraction \code{maxhomodel} of the genome (all $p$ if none qualifies).

\paragraph{Orientation and phasing.}
\code{Caller.calling} gives each cluster a state with $A \ge B$. By default
(\code{--hmm-orient state}), each bin takes its cluster's state oriented by the fitted BAF mean of
its HMM state $k = \hat z_{e,t}$: $(a_{e,t}, b_{e,t}) = (A, B)$ if $\mu_{k,2} \le \tfrac12$, else
$(B, A)$. The HMM's BAF is phased (Sec.~\ref{sec:data}), so no phasing step follows: one that
swaps whole bins would undo a true change between $(a, b)$ and $(b, a)$ at adjacent bins.

With \code{--hmm-orient bin}, CHISEL's own steps run instead. Each bin takes the state ordered by
its own BAF (\code{Caller.infer}'s rule): $(a_{e,t}, b_{e,t}) = (A, B)$ if
$\beta_{e,t} \le \tfrac12$, else $(B, A)$. \code{Caller.phase} then runs, per chromosome, a
dynamic program over the bins with two states, keep or swap the $(a, b)$ of all cells at the bin,
minimizing
$\sum_t \sum_e \bigl[\mathbb{1}(a_{e,t-1} \ne a_{e,t}) + \mathbb{1}(b_{e,t-1} \ne b_{e,t})\bigr]$
after the swaps.

\section{Interface}\label{sec:interface}

\begin{center}
\begin{tabular}{@{}lll@{}}
\toprule
 & \code{Clusterizer.kclustering} & \code{hmm.hmmclustering} \\
\midrule
data & $\mathbf{D} \in \mathbb{R}^{T \times 2N}$ (\code{format\_combo}) & same $\mathbf{D}$ \\
extra input & -- & chromosome of each row \\
hyper-parameters & \code{restarts, threshold, seed, lord, j, LB, UB} & Section~\ref{sec:params} \\
output & \code{members}: $T$ cluster labels & \code{members}: $T \times N$ state labels \\
\bottomrule
\end{tabular}
\end{center}

\section{Hyper-parameters}\label{sec:params}

\begin{center}
\begin{tabular}{@{}llll@{}}
\toprule
symbol & \code{hmmclustering} & \code{hmm\_caller.py} & default \\
\midrule
$K$ & \code{K} & \code{--hmm-states} & required \\
$\pstay$ & \code{stay} & \code{--hmm-stay} & 0.99 \\
$C_{\max}$ & \code{maxcn} & \code{--hmm-maxcn} & 6 \\
$\rho$ & \code{base} & \code{--hmm-base} & 2 \\
restarts of Section~\ref{sec:kmeans} & \code{restarts} & \code{--hmm-restarts} & 10 \\
EM iterations & \code{maxiter} & -- & 100 \\
tol & \code{tol} & -- & $10^{-2}$ \\
$\epsilon$ & \code{varfloor} & -- & $10^{-4}$ \\
$M$ & \code{maxinitpoints} (\code{BinCellHMM}) & -- & 200000 \\
\midrule
$P$ & -- & \code{-P} (CHISEL) & 4 \\
$\kappa$ & -- & \code{-A} (CHISEL) & 1 \\
$\delta$ & -- & \code{-d} (CHISEL) & 0.05 \\
\code{maxhomodel} & -- & \code{--maxhomodel} (CHISEL) & 0.2 \\
\bottomrule
\end{tabular}
\end{center}

\section{Complexity}\label{sec:complexity}

Per Baum-Welch iteration, the emissions cost $O(NTK)$ and the forward-backward recursions
$O(NTK^2)$. The posteriors and Viterbi back-pointers are held for one chromosome block at a time,
$O(N\,|\mathcal{T}_c|\,K)$ memory. The initialization computes $O(|\mathcal{P}|\,|\mathcal{G}|)$
distances.

\section{Limitations}\label{sec:limitations}

\begin{itemize}[nosep]
\item $K$ is an input; \code{kclustering} chooses its number of clusters by binary search on its
  objective.
\item One set of emission means for all cells assumes one RDR scale. A cell whose mean total copy
  number differs shifts the RDR of every state, and the HMM spends states on the shifted copies.
  Section~\ref{sec:calling} rescales each cell only after clustering.
\item The base ploidy is not identifiable from RDR and BAF: $(A, B)$ at scale $\lambda$ and
  $(2A, 2B)$ at scale $\lambda/2$ predict the same RDR and BAF. Only the penalty of \eqref{eq:bic}
  separates them~\cite{schwarz1978,mackay2003}.
\item The Gaussian BAF emission ignores the number of SNP-covering reads of each bin.
\item The BAF is taken as phased. CHISEL's \code{Combiner} orients B as each unbalanced bin's
  pseudo-bulk minor allele and flips the blocks of a balanced bin at random, so B can switch
  haplotype between adjacent bins; each switch then reads as a change between $(a, b)$ and
  $(b, a)$ states. Phase switches are not modeled.
\end{itemize}

\begin{thebibliography}{9}
\bibitem{rabiner1989} L.~R.~Rabiner. A tutorial on hidden Markov models and selected applications
  in speech recognition. \emph{Proc. IEEE} 77(2):257--286, 1989.
  \url{https://doi.org/10.1109/5.18626}
\bibitem{dempster1977} A.~P.~Dempster, N.~M.~Laird, D.~B.~Rubin. Maximum likelihood from
  incomplete data via the EM algorithm. \emph{J. R. Stat. Soc. B} 39(1):1--38, 1977.
  \url{https://doi.org/10.1111/j.2517-6161.1977.tb01600.x}
\bibitem{viterbi1967} A.~J.~Viterbi. Error bounds for convolutional codes and an asymptotically
  optimum decoding algorithm. \emph{IEEE Trans. Inf. Theory} 13(2):260--269, 1967.
  \url{https://doi.org/10.1109/TIT.1967.1054010}
\bibitem{lloyd1982} S.~P.~Lloyd. Least squares quantization in PCM. \emph{IEEE Trans. Inf.
  Theory} 28(2):129--137, 1982. \url{https://doi.org/10.1109/TIT.1982.1056489}
\bibitem{schwarz1978} G.~Schwarz. Estimating the dimension of a model. \emph{Ann. Stat.}
  6(2):461--464, 1978. \url{https://doi.org/10.1214/aos/1176344136}
\bibitem{mackay2003} D.~J.~C.~MacKay. \emph{Information Theory, Inference, and Learning
  Algorithms}. Cambridge University Press, 2003: ch.~20 and 22 (k-means, EM for mixtures),
  ch.~25 (forward-backward and Viterbi on trellises), ch.~28 (model comparison and Occam's
  razor). \url{http://www.inference.org.uk/itila/book.html}
\bibitem{isl2021} G.~James, D.~Witten, T.~Hastie, R.~Tibshirani. \emph{An Introduction to
  Statistical Learning}, 2nd ed. Springer, 2021: Section~12.4.1 (k-means), Section~6.1.3 (BIC).
  \url{https://www.statlearning.com/}
\bibitem{zaccaria2021} S.~Zaccaria, B.~J.~Raphael. Characterizing allele- and
  haplotype-specific copy numbers in single cells with CHISEL. \emph{Nat. Biotechnol.}
  39:207--214, 2021. \url{https://doi.org/10.1038/s41587-020-0661-6}
\end{thebibliography}

\end{document}
